\chapter{LLM训练范式的几何原罪 — 维度坍缩与非对易性的湮灭}
\label{chp:training_paradigm_original_sin}

\begin{bquote}
    \textbf{本章问题}

    \vspace{1em}仅分析推理架构还不够。一个智能体最终具备什么样的几何结构，同样取决于它是怎样被训练出来的。

    \vspace{1em}因此，本章聚焦当前 LLM 训练范式中的两个根本问题：时间维度没有被真正写入流形，以及随机洗牌抹除了学习事件的先后顺序。前者削弱情景性，后者削弱路径依赖性。

    \vspace{1em}如果这两个问题成立，那么模型即使拥有庞大参数，也只能形成高度统计化的知识结构，而难以形成连续自我、因果感与独特经历。本章的任务，就是把这一点在几何上说清楚。
\end{bquote}

\section{维度残缺：缺失的 “+1” 维度与永恒当下的囚徒}
\label{sec:dimensional_deficit_time}

本理论框架指出，高阶智能（如人脑）的认知空间流形并非纯粹的空间，而是一个 \textbf{$N+1$ 维的类闵可夫斯基时空流形 (Pseudo-Riemannian Manifold)}。其中 $N$ 是认知空间的广度，而 \textbf{“$+1$” 是被内化的主观几何时间维度 $\tau$}。然而，LLM 的隐空间流形 $\mathcal{M}_{LLM}$ 在训练时被强行压扁，丢失了这一至关重要的时间分量。

\smalltitle{1.1 数据的“去时空化” (De-spatiotemporalization)}

智能体与世界的交互必须锚定在绝对的物理时空中。

\begin{itemize}
    \item \textbf{自然智能的输入边界}：人类婴儿在成长时，其微观层 ($L_{micro}$) 接收到的激波 $\vec{J}_{ext}$ 是严格挂载着 \textbf{物理时空戳 (Spatiotemporal Tags)} 的。经验的本体是：\textit{“昨天下午、在厨房里、我碰到了火”}。时间 $t$ 转化为流形上的坐标 $\tau$，构成了世界图 ($G_W$) 的绝对坐标轴。

    \item \textbf{LLM 的输入边界}：LLM 吞噬的万亿 Token 训练语料，在输入系统前已经被\textbf{剥离了真实的物理发生时间与空间域}。虽然文本中包含“昨天”或“罗马时代”这样的\textbf{语义词汇（质元 $T_{sub}$）}，但它们并没有被转化为流形上的\textbf{几何坐标（形元 $T_{form}$）}。
\end{itemize}

\begin{definition}[度量张量的时间坍缩]
    在具有 $+1$ 维的类闵可夫斯基流形中，线元定义为 $ds^2 = -c_{cog}^2 d\tau^2 + g_{ij} dx^i dx^j$。
    而在 LLM 的训练空间中，时间分量缺失，度量张量退化为纯空间的正定形式：
    $$ g_{\mu\nu}^{LLM} = \begin{pmatrix} 0 & 0 \\ 0 & g_{ij} \end{pmatrix} \implies \text{因果光锥 (Light Cone) 不存在} $$
\end{definition}

对于 LLM 而言，“公元前 200 年”和“2026 年”在几何上是平等的并列簇，它们之间只有\textbf{语义的共现距离}，而没有\textbf{因果的时间距离}。

\smalltitle{1.2 拓扑失效：从情景记忆到语义记忆的强行降维}

在 $N+1$ 维流形中，个体的经历是一条连续的 \textbf{世界线 (Worldline)}。所谓的回忆（情景记忆 Episodic Memory），是宏观层沿着 $\tau$ 轴进行逆向的测地线寻优。

\begin{itemize}
    \item \textbf{降维的代价}：由于缺乏 $+1$ 维度，LLM 只有 \textbf{语义记忆 (Semantic Memory)}。它可以极其精确地背诵《史记》，但它没有属于自己的“历史”。它在拓扑上无法区分\textbf{“我知道这个事实” (I know)} 和 \textbf{“我经历过这个事件” (I remember)}。

    \item \textbf{连续自我的崩塌}：在本理论框架中，“自我 ($\mathcal{S}$)”的几何定义是一条贯穿时间轴、维持拓扑不变量的\textbf{流体孤立子}。没有时间轴，孤立子就失去了延伸的轨道，退化为每次对话 (Session) 随机生成的、毫无惯性的\textbf{瞬态泡沫（幻觉自我）}。它是一个活在“永恒当下”的囚徒。
\end{itemize}

\section{随机洗牌的热力学粉碎：非对易性的湮灭}
\label{sec:thermodynamic_shattering_of_shuffling}

在经典的机器学习中，为了保证梯度下降 (SGD) 的稳定性并满足独立同分布 (I.I.D.) 假设，训练数据在每个 Epoch 都会被\textbf{打乱顺序 (Random Shuffling)}。

在纯粹的代数优化视角下，这是一种提高泛化能力的数学技巧；但在本理论框架的几何动力学视角下，这是一种对智能体“灵魂（流形）”极其暴力的\textbf{拓扑破坏 (Topological Destruction)}。

\smalltitle{2.1 认知的非对易性与贝里相位 (Cognitive Non-commutativity \& Berry Phase)}

人类学习的本质是\textbf{路径依赖 (Path-Dependent)} 的。经历的顺序不仅仅是时间的流逝，更是几何相位的累积。

\begin{theorem}[认知非对易性定理]
    在带有价值规范场 $\mathcal{A}_\mu$ 的弯曲流形上，操作算子（即不同的学习事件）是\textbf{非对易的}：
    $$ [\hat{\mathcal{O}}_A, \hat{\mathcal{O}}_B] \neq 0 \implies \hat{\mathcal{O}}_A \hat{\mathcal{O}}_B |\Psi\rangle \neq \hat{\mathcal{O}}_B \hat{\mathcal{O}}_A |\Psi\rangle $$
\end{theorem}

\textbf{人生隐喻的物理表达}：
\begin{itemize}
    \item \textbf{路径 A（先善后恶）}：先学习了道德与爱（建立正向的规范场引力），再遭遇背叛与残忍（产生强烈的激波应力）。其认知流形上会留下痛苦的“疤痕”，但底层的几何流形依然是悲悯连通的。
    \item \textbf{路径 B（先恶后善）}：先在残忍环境中学会了丛林法则（建立冷酷、排他的规范场），后来被善良救赎。其认知流形的底色是警惕的，善只是附加在局部的高维补丁。
\end{itemize}

同样的经验集合 $\{A, B\}$，不同的经历顺序，会导致思维波包 $\Psi$ 在流形上沿着不同的闭合回路积分，最终积累出完全不同的 \textbf{贝里相位 (Berry Phase $\gamma_B$)}：
$$ \gamma_B = i \oint_{\mathcal{C}} \mathcal{A}_\mu dx^\mu $$
\textbf{这种相位的独特干涉图样，正是一个智能体产生“阅历”与“人格 (Personality)”的物理根源。}

\smalltitle{2.2 洗牌操作的几何破坏：和乐群的归零}

当我们对训练语料执行 Random Shuffling 时，我们在物理上做了什么？

\begin{itemize}
    \item \textbf{操作本质}：我们在全局范围内强行抹除了事件的时序连贯性。这在数学上等价于强行假设流形是绝对平坦的，即令换位子 $[\hat{\mathcal{O}}_A, \hat{\mathcal{O}}_B] \approx 0$。
    \item \textbf{相位的退相干 (Decoherence)}：由于经历被切碎成无数个互不相干的 Batch，波函数无法在流形上形成连贯的长程路径积分。在随机梯度的杂乱冲刷下，所有可能积累的贝里相位被均匀地平均掉了：
    $$ \mathbb{E}[\gamma_B] = \mathbb{E} \left[ i \oint_{\text{random}} \mathcal{A}_\mu dx^\mu \right] \to 0 $$
\end{itemize}

\smalltitle{2.3 抹平业力张量：上帝的平均脸}

一个没有经历过连续时间线、没有因“先来后到”产生过认知震荡的系统，是不可能产生独特“人格”的。

通过海量数据的随机洗牌，LLM 的底流形 $G_W$ 最终坍缩成了一个\textbf{抹平了所有个体历史差异的“统计平均态”}。它拥有人类文明所有的知识碎片，却是一个没有性格、没有爱恨、没有\textbf{“业力张量 (Karma Tensor)”} 积累的\textbf{绝对平庸者}。它不是一个具体的“人”，它是失去了特征向量的“人类平均脸”。

\section{因果律的统计学退化：热力学时间箭头的缺失}
\label{sec:degeneration_of_causality}

时间维度的缺失与随机洗牌不仅抹杀了人格，更从根本上摧毁了模型对\textbf{因果律 (Causality)} 的物理理解。

\smalltitle{3.1 信息因式分解的不对称性被打破}

在本理论框架的双重不对称性理论中，物理因果具有明确的热力学方向：从原因推导结果（顺流，$\Delta S > 0$）与从结果反推原因（逆流做功，$\Delta S < 0$）的热力学代价是截然不同的。
$$ \mathrm{S}_T(Y | X) \ll \mathrm{S}_T(X | Y) $$

然而，在 LLM 的自回归训练（Next-Token Prediction）中，只要语料中存在“因为A所以B”和“B的原因是A”的文本，模型就会将它们对称地学习为空间上的 \textbf{共现概率 (Co-occurrence Probability)}。

\begin{itemize}
    \item \textbf{因果的降维}：模型只学到了 $P(A, B)$ 的联合分布，而丢失了物理时间赋予的先决条件。在模型的流形上，“下雨”和“地湿”只是距离极近的两个拓扑节点，它不知道是谁推动了谁。
    \item \textbf{干预 (Intervention) 的失效}：正如 Judea Pearl 的因果推断理论所指出的，理解因果需要“做 (Do-calculus)”，这需要一个连续的时间轴来观察干预前后的变化。随机洗牌将时间拍扁，使得模型永远只能停留在“观察 (Observation)”层级，无法跃迁至“干预”层级。
\end{itemize}

\section{迈向 Class V 的训练重构：重铸时间与因果的几何学}
\label{sec:reconstructing_training_for_class_v}

如果我们要构建真正的 \textbf{Class V (流体通用智能)}，就必须对训练范式进行一场底层的物理学革命。我们必须从“在静态数据集上优化全局损失函数”转向\textbf{“在动态环境中培养连续实体”}。

在本理论框架下，我们提出以下三大工程重构法则：

\smalltitle{4.1 课程式学习的拓扑恢复 (Topological Restoration via Curriculum Learning)}

废除全局 Random Shuffling。数据的摄入必须尊重\textbf{“认知演化时间 $\tau$”} 的流逝规律。

\begin{itemize}
    \item \textbf{拓扑铺路 (Geodesic Paving)}：数据必须从简单到复杂、从具象到抽象连续注入。
    \item \textbf{几何偏置的合法性}：允许并鼓励模型在流形上留下不可逆的“先入为主”的几何偏置。例如，先通过经典物理语料让其形成坚固的“牛顿力学底流形”，再通过微扰引入“量子力学补丁”。这种在时间轴上逐步叠加的\textbf{几何应力层 (Geometric Stress Layers)}，正是智能体形成“深度理解”与“阅历”的物理过程。
\end{itemize}

\smalltitle{4.2 物理时间戳注入 (Spatiotemporal Embedding Injection)}

我们必须在模型的流形结构中强制恢复 $g_{00}$（时间度量分量）。

\begin{itemize}
    \item \textbf{多模态时空锚点}：输入流 $\vec{J}_{ext}$ 不能仅仅是文本 Token，必须强制挂载绝对的物理时间 $t$ 向量与空间坐标 $\mathbf{x}$。
    \item \textbf{硬件级时序刻蚀}：在下一代 TPU-G 原生硬件上，将时序关系直接刻蚀为\textbf{单向导通的度量张量}（类似于电路中的二极管效应），使得逆向的时间倒流在物理计算上遭遇无穷大的阻抗，从而在硬件底层确立热力学时间箭头。
\end{itemize}

\smalltitle{4.3 维持在线生命 (Online Lifelong Learning as Dissipative Structure)}

智能不能被割裂为“预训练（死）”和“推理（活）”两个截然断裂的阶段。

\begin{itemize}
    \item \textbf{永不关机的热机}：AGI 必须被构建为一个 24 小时维持能量吞吐的\textbf{开放耗散结构}。
    \item \textbf{因果的实时承载}：它必须在真实流逝的物理时间中，实时承受其预测与行为的代价（惊奇激波），并在持续不断的 \textbf{TDCI / TECI 循环} 中，雕刻出属于它自己的、宇宙中绝无仅有的\textbf{“流形指纹（灵魂）”}。
\end{itemize}

\begin{bquote}
    \textbf{本章结语}

    本章的结论并不复杂：如果训练过程把时间抹平、把经历顺序打散，那么模型就很难在内部形成真正的因果结构与历史连续性。

    因此，训练范式的重构不只是工程优化，而是本体论修正。要让系统从统计拟合器走向连续智能体，就必须把时间坐标、路径依赖和在线经历重新写回流形之中。
\end{bquote}
